\subsection{Core Language}

A \textbf{sample space} is the set of all possible outcomes of an experiment. An \textbf{event} is a subset of that sample space.

\textit{Example.} If a fair die is rolled, then
\[
S=\{1,2,3,4,5,6\}.
\]
The event ``an even number is rolled'' is
\[
E=\{2,4,6\}.
\]

When outcomes are equally likely,
\[
P(E)=\frac{\text{number of favourable outcomes}}{\text{number of possible outcomes}}.
\]

\subsection{Set Language Refresher}

Probability uses the language of sets, so it is worth recalling a few basic terms.

\textbf{Subsets.} If every element of a set $A$ also belongs to a set $B$, then $A$ is a \textbf{subset} of $B$, written
\[
A\subseteq B.
\]
In probability, an event is always a subset of the sample space.

\textbf{The universal set and complement.} The \textbf{universal set} is the set containing all objects under discussion. In probability, this role is usually played by the sample space $S$. If $A$ is a set, then its \textbf{complement} is the set of elements in the universal set that are not in $A$. It may be written as $\overline{A}$ or $A^c$.

\begin{center}
\begin{tikzpicture}[scale=0.95]
  \draw[thick] (0,0) rectangle (7,4);
  \draw[thick, fill=softivory] (3.7,2) ellipse (1.9 and 1.1);
  \node at (3.3,2.35) {$A$};
  \node at (6.25,0.35) {$S$};
  \node at (0.95,3.1) {$\overline{A}$ or $A^c$};
\end{tikzpicture}
\end{center}

\textbf{Intersection and union.} For two sets $A$ and $B$:
\begin{itemize}[leftmargin=*]
  \item $A\cap B$ is the \textbf{intersection}, meaning the elements common to both sets
  \item $A\cup B$ is the \textbf{union}, meaning the elements in $A$ or $B$ or both
\end{itemize}

\begin{center}
\begin{tikzpicture}[scale=0.95]
  \draw[thick] (0,0) rectangle (8,4.3);
  \begin{scope}
    \clip (2.9,2.1) circle (1.3);
    \fill[softivory] (5.1,2.1) circle (1.3);
  \end{scope}
  \fill[bookred!18] (2.9,2.1) circle (1.3);
  \fill[bookpurple!18] (5.1,2.1) circle (1.3);
  \draw[thick] (2.9,2.1) circle (1.3);
  \draw[thick] (5.1,2.1) circle (1.3);
  \node at (2.1,2.1) {$A$};
  \node at (5.9,2.1) {$B$};
  \node at (4,2.1) {$A\cap B$};
  \node at (7.2,0.45) {$S$};
\end{tikzpicture}
\end{center}

In a Venn diagram, the union is the whole region covered by the two circles, while the intersection is the overlap in the middle.

\subsection{Complements, Unions, and Intersections}

If $A$ is an event, then its complement means ``not $A$''. It may be written as $A^c$ or $\overline{A}$.
\[
P(A^c)=1-P(A).
\]

\textbf{Complementary events and the word ``not''.} In probability questions, the word ``not'' often signals a complement. If $E$ is an event, then the complement of $E$ is the event ``$E$ does not occur''. It is written as $E^c$ or sometimes $\overline{E}$.

\begin{center}
\begin{tikzpicture}[scale=1]
  \draw[thick] (0,0) rectangle (7,4);
  \draw[thick, fill=softivory] (3.8,2) ellipse (1.8 and 1.1);
  \node at (3.2,2.4) {$E$};
  \node at (6.2,0.4) {$S$};
  \node at (0.8,3.2) {$E^c$};
\end{tikzpicture}
\end{center}

In the diagram, the rectangle is the sample space $S$. The oval is the event $E$. Everything \emph{outside} the oval but still inside the rectangle represents $E^c$, meaning ``not $E$''.

\textit{Example.} If $E$ is ``rolling a $6$ on a fair die'', then $E^c$ is ``not rolling a $6$'', so
\[
P(E^c)=1-P(E)=1-\frac16=\frac56.
\]

For two events $A$ and $B$:
\begin{itemize}[leftmargin=*]
  \item $A\cup B$ means ``$A$ or $B$ or both''
  \item $A\cap B$ means ``$A$ and $B$''
\end{itemize}

The addition rule is
\[
P(A\cup B)=P(A)+P(B)-P(A\cap B).
\]
The subtraction avoids double-counting the outcomes common to both events.

\subsection{Mutually Exclusive Events}

Events are \textbf{mutually exclusive} if they cannot happen together. Then
\[
P(A\cap B)=0,
\]
so the addition rule simplifies to
\[
P(A\cup B)=P(A)+P(B).
\]

\textit{Warning.} Mutually exclusive does \textbf{not} mean independent. In fact, non-trivial mutually exclusive events are usually dependent.

\subsection{Multiplication Rule and Conditional Probability}

Conditional probability means the probability of an event given that another event has already happened:
\[
P(A\mid B)=\frac{P(A\cap B)}{P(B)},
\qquad P(B)>0.
\]

Rearranging gives the multiplication rule:
\[
P(A\cap B)=P(B)\,P(A\mid B)=P(A)\,P(B\mid A).
\]

This is especially useful for selection without replacement.

\subsection{Independence}

Events $A$ and $B$ are \textbf{independent} if knowing that one occurred does not change the probability of the other. Algebraically,
\[
P(A\cap B)=P(A)P(B).
\]
Equivalently, if $P(B)>0$,
\[
P(A\mid B)=P(A).
\]

\textit{Common trap.} ``Independent'' is about probabilities staying unchanged. It is not the same as ``different'' or ``mutually exclusive''.

\subsection{Counting Links}

Some probability questions are easiest when written as counting problems. Common tools include:
\begin{itemize}[leftmargin=*]
  \item permutations when order matters
  \item combinations when order does not matter
  \item complementary counting when ``at least one'' is awkward
\end{itemize}

\textit{Example.} The probability that at least one head appears in three fair coin tosses is
\[
1-P(\text{no heads})=1-\left(\frac12\right)^3=\frac78.
\]

If a question involves selecting objects, it is often worth translating the numerator and denominator into combinations. For example, if $2$ cards are chosen from a deck and we want both cards to be aces, then
\[
P(\text{both aces})=\frac{\binom42}{\binom{52}{2}}.
\]
This method is especially clean when order does not matter.

\subsection{Tree Diagrams and Staged Experiments}

Tree diagrams are useful when an experiment happens in stages. Each branch shows a probability, and the probability of a complete path is found by multiplying along the path.

\textit{Example.} Suppose a bag contains $2$ red counters and $1$ blue counter, and two counters are drawn without replacement. The tree diagram is:

\begin{center}
\begin{tikzpicture}[
  x=1cm,
  y=1cm,
  >=Latex,
  every node/.style={font=\small},
  stage/.style={circle, draw=bookpurple, thick, fill=softivory, minimum size=7mm},
  leaf/.style={rounded corners, draw=bookred, thick, fill=softivory, inner sep=4pt, align=left}
]
\node[stage] (start) at (0,0) {};

\node (r1) at (3,1.6) {Red};
\node (b1) at (3,-1.6) {Blue};

\node (rr) at (6,2.5) {Red};
\node (rb) at (6,0.7) {Blue};
\node (br) at (6,-1.6) {Red};

\node[leaf] (p_rr) at (10,2.5) {$RR:\ \frac23\cdot\frac12=\frac13$};
\node[leaf] (p_rb) at (10,0.7) {$RB:\ \frac23\cdot\frac12=\frac13$};
\node[leaf] (p_br) at (10,-1.6) {$BR:\ \frac13\cdot 1=\frac13$};

\draw[->, thick] (start) -- (r1) node[midway, above] {$\frac23$};
\draw[->, thick] (start) -- (b1) node[midway, below] {$\frac13$};

\draw[->, thick] (r1) -- (rr) node[midway, above] {$\frac12$};
\draw[->, thick] (r1) -- (rb) node[midway, below] {$\frac12$};
\draw[->, thick] (b1) -- (br) node[midway, below] {$1$};

\draw[->, thick] (rr) -- (p_rr);
\draw[->, thick] (rb) -- (p_rb);
\draw[->, thick] (br) -- (p_br);

\node[above=2pt of start] {Start};
\node[above=2pt of r1] {1st draw};
\node[above=2pt of rr] {2nd draw};
\end{tikzpicture}
\end{center}

The path probabilities add to $1$, which is a good quick check.

\subsection{Two-Way Tables}

Two-way tables are helpful when events are described by two classifications at once, such as subject choices, gender, or transport method.

\begin{center}
\begin{tabular}{lccc}
\toprule
 & $B$ & $B^c$ & Total \\
\midrule
$A$ & $n(A\cap B)$ & $n(A\cap B^c)$ & $n(A)$ \\
$A^c$ & $n(A^c\cap B)$ & $n(A^c\cap B^c)$ & $n(A^c)$ \\
Total & $n(B)$ & $n(B^c)$ & $n(S)$ \\
\bottomrule
\end{tabular}
\end{center}

Conditional probabilities then become quick to read. For example,
\[
P(A\mid B)=\frac{n(A\cap B)}{n(B)}.
\]

\subsection{Discrete Random Variables}

A \textbf{discrete random variable} takes specific numerical values, each with an associated probability. The probabilities must satisfy:
\[
0\le P(X=x)\le 1
\qquad \text{and} \qquad
\sum P(X=x)=1.
\]

These are often shown in a table.

\begin{center}
\begin{tabular}{>{\centering\arraybackslash}m{2cm} >{\centering\arraybackslash}m{3cm} >{\centering\arraybackslash}m{7cm}}
\toprule
\textbf{$x$} & \textbf{$P(X=x)$} & \textbf{Interpretation} \\
\midrule
0 & 0.2 & no successes \\
1 & 0.5 & one success \\
2 & 0.3 & two successes \\
\bottomrule
\end{tabular}
\end{center}

\textit{Bridge note.} In this booklet, the random variable takes separate values such as $0,1,2,\dots$, so we can list probabilities in a table. In continuous probability, the possible values fill intervals instead of isolated points, so the table is replaced by a density curve and sums are replaced by integrals. The central ideas stay the same: we still describe how likely values are, and we still study averages and spread.

\subsection{Expected Value and Variance}

The expected value of a discrete random variable is its weighted average:
\[
E(X)=\sum xP(X=x).
\]

The variance measures spread:
\[
\mathrm{Var}(X)=E(X^2)-[E(X)]^2,
\qquad \text{where } E(X^2)=\sum x^2P(X=x).
\]

The standard deviation is
\[
\sigma=\sqrt{\mathrm{Var}(X)}.
\]

\textit{Interpretation note.} The expected value is not necessarily an outcome that must actually occur. It is the long-run average value.

\textit{Transition remark.} This same story continues in continuous probability. Expected value still represents a long-run average, and variance still measures spread, but the formulas move from sums such as $\sum xP(X=x)$ to integrals over a probability density function. So continuous random variables are best thought of as the next version of the same ideas, not a completely different topic.

\subsection{Enrichment Note: A First Glimpse of Measure Theory}

\textit{Syllabus note.} This subsection is enrichment only. Measure theory is \textbf{outside} the HSC Mathematics Extension 1 and Extension 2 syllabuses.

At a very high level, measure theory gives a systematic way to assign ``size'' to sets. In elementary probability, that ``size'' becomes probability after normalising so that the whole sample space has total size $1$.

For a \textbf{discrete} random variable with values $x=0,1,2,\dots$, the measurement is a sum of point masses:
\[
P(X\in A)=\sum_{x\in A} P(X=x).
\]

For a single \textbf{continuous} random variable $X$, the measurement is described by a density function, and probabilities are found by integration:
\[
P(a\le X\le b)=\int_a^b f(x)\,dx.
\]

For \textbf{two variables} $X$ and $Y$, probability is measured over regions in the plane. This is why some geometric probability questions can be solved using areas:
\[
P\big((X,Y)\in R\big)=\iint_R f(x,y)\,dA.
\]

The same idea extends to higher dimensions $\mathbb{R}^n$, where probability is measured over more complicated regions. So the discrete, one-variable continuous, and area-based probability ideas in this booklet can all be viewed as early examples of one larger measuring framework.

\subsection{A Worked Variance Example}

Suppose
\[
P(X=0)=0.25,\qquad P(X=1)=0.50,\qquad P(X=2)=0.25.
\]
Then
\[
E(X)=0(0.25)+1(0.50)+2(0.25)=1.
\]
Also,
\[
E(X^2)=0^2(0.25)+1^2(0.50)+2^2(0.25)=1.5.
\]
So
\[
\mathrm{Var}(X)=1.5-1^2=0.5.
\]
This example shows that a random variable can have average value $1$ while still spreading around that centre.

\subsection{Binomial Distribution}

A random variable $X$ has a \textbf{binomial distribution} when:
\begin{itemize}[leftmargin=*]
  \item there are a fixed number $n$ of trials
  \item each trial has two outcomes, success or failure
  \item the probability of success is constant, say $p$
  \item the trials are independent
\end{itemize}

Then
\[
P(X=r)=\binom{n}{r}p^r(1-p)^{n-r}
\qquad \text{for } r=0,1,2,\dots,n.
\]

Also,
\[
E(X)=np,
\qquad
\mathrm{Var}(X)=np(1-p).
\]

\textit{Worked example.} If a biased coin has probability $0.6$ of landing heads and is tossed $4$ times, then the probability of exactly $3$ heads is
\[
\binom43(0.6)^3(0.4)=4(0.216)(0.4)=0.3456.
\]

The expected number of heads is
\[
E(X)=np=4(0.6)=2.4.
\]

\subsection{Sample Proportion as a Random Variable}

Suppose a Bernoulli trial has probability of success $p$, and it is repeated $n$ times independently. If $X$ is the number of successes, then
\[
X\sim \mathrm{Bin}(n,p).
\]

The \textbf{sample proportion} of successes is
\[
\hat p=\frac{X}{n}.
\]

This is also a random variable. Its mean is
\[
E(\hat p)=p,
\]
and its variance is
\[
\mathrm{Var}(\hat p)=\frac{p(1-p)}{n}.
\]

\textit{Interpretation note.} The sample proportion fluctuates from sample to sample, but on average it centres around the true population proportion $p$.

\subsection{Normal Approximation to the Binomial}

When $n$ is large and the binomial distribution is not too skewed, a binomial random variable can often be approximated by a normal distribution.

If
\[
X\sim \mathrm{Bin}(n,p),
\]
then a standard approximation is
\[
X \approx N\big(np,\ np(1-p)\big).
\]

Equivalently, for the sample proportion,
\[
\hat p \approx N\left(p,\ \frac{p(1-p)}{n}\right).
\]

This lets us use standard normal methods to estimate probabilities that would otherwise require many separate binomial terms.

\textit{Continuity-correction note.} When approximating a discrete variable such as $X$ by a continuous normal model, it is often better to shift endpoints by $0.5$. For example,
\[
P(55\le X\le 65)
\]
is often approximated by
\[
P(54.5\le Y\le 65.5),
\]
where $Y$ is the corresponding normal random variable.

\textit{Validity warning.} A normal approximation may be poor if the sample size is small or if the success probability is very close to $0$ or $1$. In that case the binomial distribution can be too skewed for the normal curve to be reliable.

\subsection{Common Language Triggers}

The wording of a probability question often suggests the method:
\begin{itemize}[leftmargin=*]
  \item \textbf{at least one}: usually try a complement
  \item \textbf{given that}: conditional probability
  \item \textbf{independent}: check whether probabilities stay unchanged
  \item \textbf{exactly}: count one precise case
  \item \textbf{at most}: add cases up to a bound or subtract a complement
  \item \textbf{without replacement}: probabilities usually change from stage to stage
\end{itemize}

\subsection{Method Selection Checklist}

Before calculating, ask:
\begin{itemize}[leftmargin=*]
  \item Is this a plain event probability, a conditional probability, or a random-variable question?
  \item Are the outcomes equally likely?
  \item Is replacement involved?
  \item Would a complement be quicker?
  \item Is this genuinely binomial, or do the trial conditions fail?
\end{itemize}
